\documentclass[11pt,a4paper]{article}
% Shared preamble for RuPhO English problem statements (match T1 2026 style).
\usepackage[margin=1in]{geometry}
\usepackage{amsmath,amssymb}
\usepackage{graphicx}
\usepackage{float}
\usepackage{enumitem}
\usepackage{microtype}
\usepackage{caption}
\usepackage{tabularx}
\usepackage{hyperref}

\captionsetup{font=small,labelfont=bf,skip=6pt}

\setlist[itemize]{leftmargin=1.2cm}
\setlist[enumerate]{leftmargin=1.2cm}

\newcommand{\problembox}[3]{%
  \par\addvspace{0.42em}%
  \noindent
  \setlength{\fboxsep}{7.5pt}%
  \setlength{\fboxrule}{0.95pt}%
  \fbox{%
    \begin{minipage}{\dimexpr\linewidth-2\fboxsep-2\fboxrule\relax}
    \begin{tabularx}{\linewidth}{@{}>{\raggedright\arraybackslash\bfseries}p{2.1em} @{\hspace{0.9em}} >{\raggedright\arraybackslash}X @{\hspace{0.9em}} >{\raggedleft\arraybackslash\bfseries}p{2.4em}@{}}
    #1 & #2 & #3
    \end{tabularx}
    \end{minipage}%
  }%
  \par\addvspace{0.32em}%
}

\title{\textbf{Conductors in a magnetic field --- Solution}}
\author{\normalsize X24 (2024) — English translation}
\date{}
\begin{document}
\maketitle
\problembox{A1}{Let the moment of time $t_0=0$ the load be at the origin of coordinates, and the projection of its velocity onto the axis $x$ is equal to $v_0$. Determine the dependences of the coordinate $x(t)$ and speed $v_x(t)$ of the load on time $t$. Express your answer using $v_0$, $\gamma$, $\omega_0$ and $t$.\\
Let's write down the equation of motion of the load: $$m\ddot{x}=-\beta\dot{x}-kx\Rightarrow \ddot{x}+\cfrac{\beta}{m}\dot{x}+\cfrac{k}{m}x=0{.} $$   : $$\ddot{x}+2\gamma\dot{x}+\omega^2_0x=0{.} $$   in  : $$x=Re\left\{Ae^{\lambda t}\right\}{,} $$where $A\neq 0$ and $\lambda$ –   . : $$A\left(\lambda^2+2\gamma\lambda+\omega^2_0\right)=0\Rightarrow{\lambda=-\gamma\pm{i\sqrt{\omega^2_0-\gamma^2}}}{.} $$   mm ,    $\lambda$: $$x(t)=e^{-\gamma t}Re\left\{A_1e^{i\sqrt{\omega^2_0-\gamma^2}t}+A_2e^{-i\sqrt{\omega^2_0-\gamma^2}t}\right\}=Ce^{-\gamma t}\sin\left(\sqrt{\omega^2_0-\gamma^2}t+\varphi_0\right){,} $$where $C$ and $\ varphi_0$ –  . : $$v_x(t)=Ce^{-\gamma t}\left(\sqrt{\omega^2_0-\gamma^2}\cos\left(\sqrt{\omega^2_0-\gamma^2}t+\varphi_0\right)-\gamma\sin\left(\sqrt{\omega^2_0-\gamma^2}t+\varphi_0\right)\right){.} $$   $t=0$ : $$\begin{cases} x(0)=C\sin\varphi_0\\ v_x(0)=C\left(\sqrt{\omega^2_0-\gamma^2}\cos\varphi_0-\gamma\sin\varphi_0\right) \end{cases} $$   : $$\begin{cases} x(0)=0\\ v_x(0)=v_0 \end{cases} \Rightarrow \begin{cases} \varphi_0=0\\ C=\cfrac{v_0}{\sqrt{\omega^2_0-\gamma^2}} \end{cases} $$Final we get:\\
Answer: $$x(t)=\cfrac{v_0}{\sqrt{\omega^2_0-\gamma^2}}e^{-\gamma t}\sin\left(\sqrt{\omega^2_0-\gamma^2}t\right){.} $$}{0.60}
\problembox{A2}{Obtain an exact expression for $Q$. Express your answer in terms of $\omega_0$ and $\gamma$.\\
Let's rewrite the expression for the quality factor in the following form: $$Q=\cfrac{2\pi}{1-\left(\cfrac{v_1}{v_0}\right)^2}{,} $$where $v_1$ –    for     s    .   $v_1$    $T$, : $$T=\cfrac{2\pi}{\sqrt{\omega^2_0-\gamma^2}}{.} $$: $$\cfrac{v_1}{v_0}=e^{-\gamma T}=e^{-2\pi\gamma/\sqrt{\omega^2_0-\gamma^2}}{,} $$from where we get:\\
Answer: $$Q=\cfrac{2\pi}{1-e^{-4\pi\gamma/\sqrt{\omega^2_0-\gamma^2}}}{.} $$}{0.40}
\problembox{A3}{Derive an approximate expression for the quality factor $Q$ at weak attenuation ($\gamma\ll\omega_0$). Express your answer using $m$, $k$ and $\beta$.\\
The expansion of the denominator at $\gamma\ll{\omega_0}$ is as follows: $$1-e^{-4\pi\gamma/\sqrt{\omega^2_0-\gamma^2}}\approx 1-e^{-4\pi\gamma/\omega_0}\approx 1-\left(1-\cfrac{4\pi\gamma}{\omega_0}\right)=\cfrac{4\pi\gamma}{\omega_0}{,} $$   for  for   $Q$: $$Q\approx\cfrac{\omega_0}{2\gamma}{.} $$ $\omega_0$ and $\gamma$, we obtain:\\
Answer: $$Q\approx\cfrac{\sqrt{mk}}{\beta(H)}{.} $$}{0.20}
\problembox{B1}{The deviation $x$ of the load from the position depends on time $t$ as follows: $$x(t)=A\sin\left(\Omega t+\varphi_0\right) $$ $A$ and $\varphi_0$.    $A_0$, $\Omega$, $\omega_0$ and $\gamma$.\\
The equation of motion is as follows: $$m\ddot{x}=k(A_0\sin\Omega t-x)-\beta\dot{x}{,} $$: $$\ddot{x}+2\gamma\dot{x}+\omega^2_0\Delta{x}=\omega^2_0A_0\sin\Omega t{.} $$   in : $$x=Re\left\{\hat{A}e^{i\Omega t}\right\}{,} $$where $\hat{A}\neq 0$ –   . : $$\hat{A}\left((\omega^2_0-\Omega^2)+2i\Omega\gamma\right)=\omega^2_0A_0e^{-i\pi/2}{.} $$ $A$: $$\hat{A}=\cfrac{\omega^2_0A_0\left((\omega^2_0-\Omega^2)-2i\Omega\gamma\right)e^{-i\pi/2}}{\left((\omega^2_0-\Omega^2)^2+4\gamma^2\Omega^2\right)}=\cfrac{A_0\omega^2_0e^{-i(\pi/2-\varphi)}}{\sqrt{(\omega^2_0-\Omega^2)^2+4\gamma^2\omega^2_0}}{,} $$where $\varphi$ –    $z=(\omega^2_0-\Omega^2) -2i\Omega\gamma$.   h, : $$\Delta{x}(t)=\cfrac{A_0\omega^2_0}{\sqrt{(\omega^2_0-\Omega^2)^2+4\gamma^2\Omega^2}}\sin\left(\Omega t+\varphi\right){.} $$ : $$A=\cfrac{A_0\omega^2_0}{\sqrt{(\omega^2_0-\Omega^2)^2+4\gamma^2\Omega^2}}\qquad \varphi_0=\varphi{.} $$  $\varphi$   : $$\varphi=\begin{cases} -\arctan\cfrac{2\gamma\Omega}{\omega^2_0-\Omega^2}\quad\text{for}\quad \Omega{<}\omega_0\\ -\cfrac{\pi}{2}\quad\text{for}\quad \Omega=\omega_0\\ -\pi-\arctan\cfrac{2\gamma\Omega}{\omega^2_0-\Omega^2}\quad\text{for}\quad \Omega{>}\omega_0 \end{cases} $$So way:\\
Answer: $$A=\cfrac{A_0\omega^2_0}{\sqrt{(\omega^2_0-\Omega^2)^2+4\gamma^2\Omega^2}}{.} $$$$\varphi_0=\begin{cases} -\arctan\cfrac{2\gamma\Omega}{\omega^2_0-\Omega^2}\quad\text{for}\quad \Omega{<}\omega_0\\ -\cfrac{\pi}{2}\quad\text{for}\quad \Omega=\omega_0\\ -\pi-\arctan\cfrac{2\gamma\Omega}{\omega^2_0-\Omega^2}\quad\text{for}\quad \Omega{>}\omega_0 \end{cases} $$}{0.60}
\problembox{B2}{Obtain exact expressions for the resonant cyclic frequency $\Omega_{}$ and the corresponding oscillation amplitude $A_{}$. Express your answers using $\omega_0$, $\gamma$ and $A_0$. Consider that $\gamma\sqrt{2}<\omega_0$.\\
Differentiating the denominator with respect to $\Omega^2$, we obtain: $$-2(\omega^2_0-\Omega^2)+4\gamma^2=0{,} $$from where:\\
Substituting $\Omega_\text{}$ into the expression for $A$, we find:\\
Answer: $$\Omega_\text{}=\sqrt{\omega^2_0-2\gamma^2}{.} $$}{0.30}
\problembox{B3}{Derive approximate expressions for $\Omega_{}$, $A_{}$ and $\Delta{\omega}$ under weak damping ($\gamma\ll{\omega_0}$). Express your answers using $A_0$, $\omega_0$ and $\gamma$.\\
Simplified expressions for $\Omega_\text{}$ and $A_\text{}$ are obtained trivially and take the following form:\\
Consider the cyclic frequency $\Omega=\Omega_\text{}+\Delta\Omega$, where $\Delta\Omega\ll\Omega_\text{}$. The value of $\Omega_\text{}$ can be considered equal to $\omega_0$, since: $$\Omega_\text{}=\sqrt{\omega^2_0-2\gamma^2}\approx \omega_0-\cfrac{\gamma^2}{\omega_0}{.} $$ $\Omega_\text{res}$ from $\omega_0$      .  for   : $$\left(\omega^2_0-\Omega^2\right)^2+4\gamma^2\Omega^2\approx \left(\omega^2_0-\left(\omega_0+\Delta\Omega\right)\right)^2+4\gamma^2\left(\omega_0+\Delta\Omega\right)^2\approx 4\omega^2_0\Delta\Omega^2+4\gamma^2\omega^2_0{.} $$ $\Delta{\Omega }$ ,       . : $$4\Delta{\Omega}^2\omega^2_0+4\gamma^2\omega^2_0=8\gamma^2\omega^2_0\Rightarrow \Delta{\Omega}_{1.2}=\pm\gamma{.} $$ $\Delta\omega=\Delta\Omega_1-\Delta\Omega_2$, we get:\\
Answer: $$\Omega_\text{}\approx{\omega_0}\qquad A_\text{}\approx{\cfrac{A_0\omega_0}{2\gamma}}{.} $$}{0.30}
\problembox{C1}{Find the induction $B_x$ of the magnetic field of the ring on its axis at the point with coordinate $x$. Express your answer in terms of $x$, $R$, $I$ and the magnetic constant $\mu_0$.\\
From the Biot-Savart-Laplace law: $$d\vec{B}=\cfrac{\mu_0}{4\pi}\cfrac{\left[\vec{r}\times d\vec{r}\right]}{r^3}{,} $$where $\vec{r}$ – –     s  $x$.    : $$r=\sqrt{R^2+x^2}{.} $$   from       cm: $$\oint_L\left[\vec{r}\times d\vec{r}\right]=2\vec{S}=2\pi{R}^2\vec{e}_x{,} $$from:\\
Answer: $$B_x(x)=\cfrac{\mu_0IR^2}{2\left(R^2+x^2\right)^{3/2}}{.} $$}{0.30}
\problembox{C2}{Determine the magnetic moment $\vec{m}$ of the disk. Express your answer in terms of $\vec{e}_x$, $r_0$, $h$, $\rho$ and $\dot{B}$.\\
From Faraday's law of electromagnetic induction we obtain: $$\mathcal{E}_\text{}=-\cfrac{d\Phi}{dt}=-\pi{r}^2\dot{B}=2\pi rE_\text{}{,} $$: $$E_\text{}=-\cfrac{\dot{B}r}{2}{.} $$  \Omega in  : $$\vec{j}=\cfrac{\vec{E}}{\rho}{,} $$: $$j(r)=-\cfrac{\dot{B}r}{2\rho}{.} $$    $h$ and  $dr$ : $$dm_x=\pi{r}^2dI=\pi{r}^2jhdr=-\cfrac{\pi\dot{B}h}{2\rho}r^3dr{,} $$: $$m_x=-\cfrac{\pi\dot{B}h}{2\rho}\int\limits_0^{r_0}r^3dr{.} $$Integrating, we find:\\
Answer: $$\vec{m}=-\vec{e}_x\cdot{\cfrac{\pi{r}^4_0h\dot{B}}{8\rho}}{.} $$}{1.00}
\problembox{C3}{Determine the magnetic moment $\vec{m}$ of the ball. Express your answer in terms of $\vec{e}_x$, $R_0$, $\rho$ and $\dot{B}$.\\
Let's use spherical coordinates (that is, we will count the angle $\theta$ from the positive direction of the axis $x$). Then for $r_0$ and $h$ we have: $$r_0=R_0\sin\theta\qquad h=d(R_0(1-\cos\theta))=R_0\sin\theta d\theta{.} $$ for     : $$dm_x=-\cfrac{\pi{R}^5_0\dot{B}}{8\rho}\sin^5\theta d\theta{,} $$: $$m_x=-\cfrac{\pi{R}^5_0\dot{B}}{8\rho}\int\limits_{0}^{\pi}\sin^5\theta d\theta{.} $$   s   $x=\cos\theta$: $$\int\limits_{0}^{\pi}\sin^5\theta d\theta=\int\limits_{-1}^{1}(1-x^2)^2dx=\int\limits_{-1}^{1}(1-2x^2+x^4)dx=\left(x-\cfrac{2x^3}{3}+\cfrac{x^5}{5}\right)\Biggl|_{-1}^{1}=\cfrac{16}{15}{,} $$from:\\
Answer: $$\vec{m}=-\vec{e}_x\cdot\cfrac{2\pi{R}^5_0\dot{B}}{15\rho}{.} $$}{0.50}
\problembox{C4}{Obtain the time derivative of the magnetic field induction of the ring at the center of the ball $dB_x/dt$, which is equivalent to the value $\dot{B}$. Express your answer in terms of $v$, $I$, $R$, $x$ and the magnetic constant $\mu_0$.\\
When the ball moves, the magnetic field induction $B_x(x)=B_x(x(t))$, i.e., is a complex function of time, so we have: $$\dot{B_x}=\cfrac{dB_x}{dt}=\cfrac{dB_x}{dx}\cfrac{dx}{dt}=v\cfrac{dB_x}{dx} $$  $dB_x/dx$: $$\cfrac{dB_x}{dx}=-\cfrac{3\mu_0IR^2x}{2(R^2+x^2)^{\frac{5}{2}}}{,} $$from where:\\
Answer: $$\dot{B}=-\cfrac{3\mu_0IR^2xv}{2(R^2+x^2)^{\frac{5}{2}}}{.} $$}{0.40}
\problembox{C5}{Find the proportionality factor $\beta(x)$. Express your answer in terms of $I$, $R$, $x$, $R_0$, $\rho$ and the magnetic constant $\mu_0$.\\
Since the ball moves along the $x$ axis: $$\vec{F}=m_x\cdot\cfrac{\partial\vec{B}}{\partial x}=\vec{e}_x\cdot m_x\cfrac{dB_x}{dx}{.} $$   $m_x$ : $$m_x=-\cfrac{2\pi{R}^5_0v}{15\rho}\cdot\cfrac{dB_x}{dx}{,} $$: $$F_x=-\cfrac{2\pi{R}^5_0}{15\rho}\left(\cfrac{dB_x}{dx}\right)^2v{.} $$ $dB_x/dx$, we find:\\
Answer: $$\beta(x)=\cfrac{3\pi\mu^2_0I^2R^4R^5_0x^2}{10\rho(R^2+x^2)^5}{.} $$}{0.50}
\problembox{C6}{Determine the resistivity $\rho$ of the ball used in the first experiment. Express your answer in terms of $m$, $k$, $R_0$, $R$, $H$, $I$ and the magnetic constant $\mu_0$.\\
The first peak is reached at $\Delta{x}\approx{47~\text{.}}$, and the 12th at $\Delta{x}\approx{6~\text{.}}$. Hence: $$\cfrac{\gamma}{\sqrt{\omega^2_0-\gamma^2}}\approx{\cfrac{\ln\left(\cfrac{47}{6}\right)}{11\cdot{2\pi}}}\approx{0.0297}\Rightarrow \cfrac{\gamma}{\omega_0}\approx 0{.}03\ll{1}{.} $$  : $$\cfrac{\omega_0}{2\gamma}=\cfrac{\sqrt{mk}}{\beta(H)}{,} $$: $$\beta(H)\approx{\cfrac{2\gamma\sqrt{mk}}{\omega_0}}{,} $$and finally\\
Answer: $$\rho=15{.}7\cdot\cfrac{\mu^2_0I^2R^5_0R^4H^2}{\sqrt{mk}(R^2+H^2)^5} $$}{0.80}
\problembox{C7}{Determine the resistivity $\rho$ of the ball used in the second experiment. Express your answer in terms of $m$, $k$, $R_0$, $R$, $H$, $I$ and the magnetic constant $\mu_0$.\\
From the expression for the resonant amplitude we find: $$Q=\cfrac{A_\text{}}{A_0}=\cfrac{\sqrt{mk}}{\beta(H)}\approx 25 $$from where finally:\\
Answer: $$\rho=23{.}6\cfrac{\mu^2_0I^2R^5_0R^4H^2}{\sqrt{mk}(R^2+H^2)^5} $$}{0.70}
\problembox{D1}{Determine the induction $B_z$ of the magnetic field of the solenoid, as well as its derivative $dB_z/dz$ at the point with coordinate $z$. Express your answer using $\mu_0$, $n$, $I$, $R$ and $z$.\\
From the solid angle theorem for the magnetic field we have: $$B_z=\cfrac{\mu_0i\Omega_\text{}}{4\pi}{,} $$where $i=In$ –    ,  $\Omega_\text{side}$ –  ,      .  ,      ,   ,        s  $z$,  : $$\Omega_\text{}=2\pi(1-\cos\alpha)=2\pi\left(1-\cfrac{z}{\sqrt{z^2+R^2}}\right){.} $$Thus:\\
Differentiating, we find:\\
Answer: $$B_z=\cfrac{\mu_0nI}{2}\left(1-\cfrac{z}{\sqrt{z^2+R^2}}\right){.} $$}{0.60}
\problembox{D2}{Determine the linear current density $i$ on the surface of the cylinder at the point with coordinate $z$. Express your answer using $\mu_0$, $x$ and $dB_z(z)/dz$.\\
Since the magnetic field induction outside the cylinder can be considered equal to the magnetic field induction of the solenoid, we have: $$B_{z(in)}=B_z(z-x)\qquad B_{z(out)}=B(z){.} $$    for    : $$(B_{z(in)}-B_{z(out)})\approx -x\cfrac{dB_z}{dz}=\mu_0ix{,} $$from where:\\
Answer: $$i(z)=-\cfrac{x}{\mu_0}\cfrac{dB_z}{dz}{.} $$}{1.00}
\problembox{D3}{Determine the force $F_x$ acting on the cylinder from the magnetic field of the solenoid. Express your answer in terms of $\mu_0$, $r$, $R$, $n$, $I$ and $x$.\\
The magnetic moment of an infinitesimal section of a cylinder is: $$dm_z=i(z)\pi r^2dz{.} $$  on   $dF_x$ : $$dF_x=dm_z\cfrac{dB_z}{dz}{.} $$  for $i$, : $$dF_x=-\cfrac{\pi{r}^2x}{\mu_0}\left(\cfrac{dB_z}{dz}\right)^2dz{.} $$  for $dB_z/dz$: $$dF_x=-\cfrac{\mu_0\pi{r}^2n^2I^2xR^4dz}{4(R^2+z^2)^3}{,} $$ s      from  : $$F_x\approx -\cfrac{\mu_0\pi{r}^2n^2I^2R^4x}{4}\int\limits_{-\infty}^{\infty}\cfrac{dz}{(R^2+z^2)^3}{.} $$   $z=R \operatorname{tg}\varphi$ and : $$F_x=\cfrac{\mu_0\pi{r}^2n^2I^2R^4x}{4}\int\limits_{-\infty}^{\infty}\cfrac{dz}{(R^2+z^2)^3}=\cfrac{\mu_0\pi{r}^2n^2I^2x}{4R}\int\limits_{-\pi/2}^{\pi/2}\cos^4\varphi d\varphi{.} $$  : $$\int\limits_{-\pi/2}^{\pi/2}\cos^4\varphi d\varphi=\int\limits_{-\pi/2}^{\pi/2}\left(\cfrac{1+\cos 2\varphi}{2}\right)^2d\varphi=\int\limits_{-\pi/2}^{\pi/2}\left(\cfrac{1}{4}+\cfrac{\cos 2\varphi}{2}+\cfrac{\cos^22\varphi}{4}\right)d\varphi= $$$$=\int\limits_{-\pi/2}^{\pi/2}\left(\cfrac{1}{4}+\cfrac{\cos 2\varphi}{2}+\cfrac{1}{4}\left(\cfrac{1+\cos 4\varphi}{2}\right)\right)d\varphi=\int\limits_{-\pi/2}^{\pi/2}\left(\cfrac{3}{8}+\cfrac{\cos 2\varphi}{2}+\cfrac{\cos 4\varphi}{8}\right)d\varphi= $$$$=\left(\cfrac{3\varphi}{8}+\cfrac{\sin 2\varphi}{4}+\cfrac{\sin 4\varphi}{32}\right)\Biggl|_{-\pi/2}^{\pi/2}=\cfrac{3\pi}{8}{.} $$So way:\\
Answer: $$F_x=-\cfrac{3\pi^2\mu_0\pi r^2n^2I^2}{32R}x{.} $$}{1.50}
\problembox{D4}{Obtain the dependence of the displacement of the rod $x$ on time $t$. Express your answer in terms of $\mu_0$, $r$, $R$, $n$, $I$ and $m$.\\
Let us write the equation of motion of the cylinder: $$m\ddot{x}=-\cfrac{3\mu_0\pi^2 r^2n^2I^2}{32R}x\Rightarrow \omega^2_0=\cfrac{3\mu_0\pi^2 r^2n^2I^2}{32mR}\Rightarrow x(t)=A\sin\omega_0t+B\cos\omega_0t{.} $$  $A$ and $B$   : $$ \begin{cases} x(0)=B=0\\ v_x(0)=\omega_0A=v_0 \end{cases} \Rightarrow A=\cfrac{v_0}{\omega_0} $$Thus:\\
Answer: $$x(t)=v_0\sqrt{\cfrac{32mR}{3\mu_0\pi^2{r}^2n^2I^2}}\sin\sqrt{\cfrac{3\mu_0\pi^2 r^2n^2I^2}{32mR}}t{.} $$}{0.30}
\vspace{1em}
\noindent\textit{Source:} \url{https://pho.rs/p/2912/s}
\end{document}
